% ===== 导言区：宏包、环境与宏定义（沿用线性规划笔记的框架） =====
\usepackage{amsmath,amssymb,amsthm,mathtools}
\usepackage{geometry}
\geometry{a4paper,margin=2.5cm}
\usepackage{graphicx}
\usepackage{enumitem}
\usepackage{xcolor}
\usepackage{tikz}
\usepackage{booktabs}
\usepackage{longtable}
\usepackage{environ}
\usepackage[colorlinks=true,allcolors=black]{hyperref}
\usetikzlibrary{arrows.meta,positioning,calc}

% ---------- 是否显示练习解答（true/false） ----------
\newif\ifshowsolutions
\showsolutionstrue

% ---------- 定理类环境（按章编号，共用计数器） ----------
\theoremstyle{definition}  % 正文直立，中文不宜斜体
\newtheorem{theorem}{定理}[chapter]
\newtheorem{proposition}[theorem]{命题}
\newtheorem{lemma}[theorem]{引理}
\newtheorem{corollary}[theorem]{推论}
\newtheorem{definition}[theorem]{定义}
\newtheorem{example}[theorem]{例}
\theoremstyle{remark}
\newtheorem{remark}[theorem]{注}

\renewcommand{\proofname}{证明}

% ---------- 练习与解答 ----------
% 练习按章编号；可选参数写来源，如 \begin{exercise}[原讲义 B.1 第 2 题]
\newcounter{exercise}[chapter]
\renewcommand{\theexercise}{\thechapter.\arabic{exercise}}
\newenvironment{exercise}[1][]{%
  \refstepcounter{exercise}%
  \par\medskip\noindent\textbf{练习 \theexercise}%
  \if\relax\detokenize{#1}\relax\else（#1）\fi\quad
}{\par\medskip}
\NewEnviron{solution}{%
  \ifshowsolutions
    \par\smallskip\noindent\textbf{解答}\quad\BODY\par\medskip
  \fi
}

% ---------- 待填充占位 ----------
\newcommand{\todo}[1]{{\color{red}\textbf{【待补充】}#1}}

% ---------- 常用数学宏 ----------
\newcommand{\R}{\mathbb{R}}
\newcommand{\Q}{\mathbb{Q}}
\newcommand{\Z}{\mathbb{Z}}
\newcommand{\N}{\mathbb{N}}
\newcommand{\Rbar}{\overline{\R}}            % R ∪ {+∞}
\newcommand{\X}{\mathcal{X}}
\newcommand{\Y}{\mathcal{Y}}
\newcommand{\Xs}{\mathcal{X}^{*}}
\newcommand{\Ys}{\mathcal{Y}^{*}}
\newcommand{\inner}[2]{\langle #1, #2\rangle}
\newcommand{\pair}[2]{\langle #1 \mid #2\rangle}  % 对偶配对 <x^*|x>
\newcommand{\norm}[1]{\lVert #1\rVert}
\newcommand{\abs}[1]{\lvert #1\rvert}
\newcommand{\Gammaz}{\Gamma_{0}}            % 正常下半连续凸函数类
\DeclareMathOperator{\conv}{conv}
\newcommand{\clconv}{\overline{\conv}}       % 闭凸包
\DeclareMathOperator{\dom}{dom}
\DeclareMathOperator{\epi}{epi}
\DeclareMathOperator{\interior}{int}
\DeclareMathOperator{\cl}{cl}
\DeclareMathOperator{\ri}{ri}
\DeclareMathOperator{\aff}{aff}
\DeclareMathOperator{\cone}{cone}
\DeclareMathOperator{\spanop}{span}
\DeclareMathOperator{\cont}{cont}
\DeclareMathOperator{\lev}{lev}
\DeclareMathOperator{\Nor}{Nor}
\DeclareMathOperator{\prox}{prox}
\DeclareMathOperator{\Lip}{Lip}
\DeclareMathOperator{\Id}{Id}
\DeclareMathOperator{\supp}{supp}
\DeclareMathOperator*{\argmin}{arg\,min}
\DeclareMathOperator*{\argmax}{arg\,max}
\newcommand{\T}{\top}

% ---------- 复习辅助 ----------
% 标题后加 \opt 表示“了解即可”的小节
\newcommand{\opt}{\textsuperscript{$\ast$}}
% 每章末尾的要点清单
\newenvironment{keypoints}{%
  \par\bigskip\noindent\textbf{本章要点}%
  \begin{itemize}[leftmargin=2em,itemsep=1pt,topsep=3pt]}{\end{itemize}\par\medskip}

% ---------- 按约定引入各章正文与练习 ----------
% 正文：chapters/chNN.tex（自带 \chapter）；练习：exercises/chNN.tex（不含 \chapter）
\newcommand{\includechapter}[1]{%
  \IfFileExists{chapters/ch#1.tex}%
    {\input{chapters/ch#1}}%
    {\chapter{第 \number#1\relax{} 章（待填充）}\todo{尚未完成。}}%
  \IfFileExists{exercises/ch#1.tex}%
    {\section{练习}\input{exercises/ch#1}}%
    {}%
}
